% GENERATED FILE. Edit canonical lessons, then run npm run generate:books.
% canonical-source-hash: fnv1a64:6112d91fb4395851
% canonical-lessons: MR-C313-aphva, MR-C313-khalbal, MR-C313-ghosna, MR-C313-thalak, MR-C313-pukarne

\chapter{Rumour and Sensation}
\label{ch:mr-rumour-and-sensation}
\hlchaptermodality{313}

\begin{chapteropening}
\textbf{By the end of this chapter:} \emph{I can talk about rumours, announcements and a stir in town.}

Complete the last lesson of chapter 313: I can talk about rumours, announcements and a stir in town.
\end{chapteropening}

\section[\texorpdfstring{\mr{अफवा} (aphvā)}{अफवा (aphvā)}]{\mr{अफवा} (aphvā) — a rumour}

\label{lesson:MR-C313-aphva}



\begin{quote}
Before the new one: say the Marathi for a newsreader, then the Marathi for to broadcast.
\end{quote}

\subsection*{You'll want to know: \mr{अफवा}}
\textbf{\mr{अफवा}} — \emph{aphvā} — \textquotedblleft{}a rumour\textquotedblright{}.

\begin{quote}
\mr{गावात} \mr{भूकंपाची} \mr{खोटी} \mr{अफवा} \mr{पसरली}.
\end{quote}

\begin{quote}
\emph{gāvāt bhūkampācī khoṭī aphvā pasarlī.}
\end{quote}

\begin{quote}
A false rumour of an earthquake spread through the village.
\end{quote}

Feminine.

\mr{अफवा} \mr{पसरणे}, a rumour spreads; plural \mr{अफवा}.

\begin{cousinweb}
\textbf{\mr{अफवा}} (\emph{aphvā}) is a loan from Persian \emph{afvāh}, usually traced to Arabic \emph{afwāh}, 'mouths': talk that goes from mouth to mouth.
\end{cousinweb}

\begin{grammarlens}[title={What you've built}]
One more word for telling what happened.
\end{grammarlens}

\subsection*{Your turn}
\begin{itemize}
\raggedright
  \item \emph{Say it:} \emph{aphvā}
  \item \emph{Say it:} \emph{aphvā}, once more
  \item \emph{Say it:} \emph{gāvāt bhūkampācī khoṭī aphvā pasarlī.} — A false rumour of an earthquake spread through the village.
  \item \emph{Recall:} say the Marathi for a newsreader, then the Marathi for to broadcast, then say \emph{aphvā} again
\end{itemize}

\begin{tcolorbox}[breakable,colback=teal!4,colframe=teal!35!black,title={Before you move on}]
What does \mr{अफवा} mean? (\textquotedblleft{}A rumour\textquotedblright{}.) Say it once more.
\end{tcolorbox}
\section[\texorpdfstring{\mr{खळबळ} (khaḷbaḷ)}{खळबळ (khaḷbaḷ)}]{\mr{खळबळ} (khaḷbaḷ) — a sensation; a stir}

\label{lesson:MR-C313-khalbal}



\begin{quote}
Before the new one: say the Marathi for to broadcast, then the Marathi for a rumour.
\end{quote}

\subsection*{You'll want to know: \mr{खळबळ}}
\textbf{\mr{खळबळ}} — \emph{khaḷbaḷ} — \textquotedblleft{}a sensation; a stir\textquotedblright{}.

\begin{quote}
\mr{त्या} \mr{बातमीने} \mr{शहरात} \mr{खळबळ} \mr{उडाली}.
\end{quote}

\begin{quote}
\emph{tyā bātmīne śahrāt khaḷbaḷ uḍālī.}
\end{quote}

\begin{quote}
That news caused a sensation in the city.
\end{quote}

Feminine.

\mr{खळबळ} \mr{उडणे}, a stir breaks out; \mr{खळबळजनक}, sensational.

\begin{cousinweb}
\textbf{\mr{खळबळ}} (\emph{khaḷbaḷ}) is inherited: CDIAL lists Marathi \emph{khaḷbaḷṇe}, 'to agitate', under Sanskrit \emph{khalati}, 'shakes' (CDIAL 3837); the noun names the stir.
\end{cousinweb}

\begin{grammarlens}[title={What you've built}]
One more word for telling what happened.
\end{grammarlens}

\subsection*{Your turn}
\begin{itemize}
\raggedright
  \item \emph{Say it:} \emph{khaḷbaḷ}
  \item \emph{Say it:} \emph{khaḷbaḷ}, once more
  \item \emph{Say it:} \emph{tyā bātmīne śahrāt khaḷbaḷ uḍālī.} — That news caused a sensation in the city.
  \item \emph{Recall:} say the Marathi for to broadcast, then the Marathi for a rumour, then say \emph{khaḷbaḷ} again
\end{itemize}

\begin{tcolorbox}[breakable,colback=teal!4,colframe=teal!35!black,title={Before you move on}]
What does \mr{खळबळ} mean? (\textquotedblleft{}A sensation; a stir\textquotedblright{}.) Say it once more.
\end{tcolorbox}
\section[\texorpdfstring{\mr{घोषणा} (ghoṣṇā)}{घोषणा (ghoṣṇā)}]{\mr{घोषणा} (ghoṣṇā) — an announcement; a slogan}

\label{lesson:MR-C313-ghosna}



\begin{quote}
Before the new one: say the Marathi for a rumour, then the Marathi for a sensation.
\end{quote}

\subsection*{You'll want to know: \mr{घोषणा}}
\textbf{\mr{घोषणा}} — \emph{ghoṣṇā} — \textquotedblleft{}an announcement; a slogan\textquotedblright{}.

\begin{quote}
\mr{स्टेशनवर} \mr{गाडी} \mr{उशिरा} \mr{येण्याची} \mr{घोषणा} \mr{झाली}.
\end{quote}

\begin{quote}
\emph{sṭeśanvar gāḍī uśirā yeṇyācī ghoṣṇā jhālī.}
\end{quote}

\begin{quote}
At the station it was announced that the train would be late.
\end{quote}

Feminine.

\mr{घोषणा} \mr{करणे}, to announce.

\begin{cousinweb}
\textbf{\mr{घोषणा}} (\emph{ghoṣṇā}) is a tatsama, taken from Sanskrit \emph{ghoṣaṇā}, 'a proclaiming aloud' (Monier-Williams).
\end{cousinweb}

\begin{grammarlens}[title={What you've built}]
One more word for telling what happened.
\end{grammarlens}

\subsection*{Your turn}
\begin{itemize}
\raggedright
  \item \emph{Say it:} \emph{ghoṣṇā}
  \item \emph{Say it:} \emph{ghoṣṇā}, once more
  \item \emph{Say it:} \emph{sṭeśanvar gāḍī uśirā yeṇyācī ghoṣṇā jhālī.} — At the station it was announced that the train would be late.
  \item \emph{Recall:} say the Marathi for a rumour, then the Marathi for a sensation, then say \emph{ghoṣṇā} again
\end{itemize}

\begin{tcolorbox}[breakable,colback=teal!4,colframe=teal!35!black,title={Before you move on}]
What does \mr{घोषणा} mean? (\textquotedblleft{}An announcement; a slogan\textquotedblright{}.) Say it once more.
\end{tcolorbox}
\section[\texorpdfstring{\mr{ठळक} (ṭhaḷak)}{ठळक (ṭhaḷak)}]{\mr{ठळक} (ṭhaḷak) — prominent; bold}

\label{lesson:MR-C313-thalak}



\begin{quote}
Before the new one: say the Marathi for a sensation, then the Marathi for an announcement.
\end{quote}

\subsection*{You'll want to know: \mr{ठळक}}
\textbf{\mr{ठळक}} — \emph{ṭhaḷak} — \textquotedblleft{}prominent; bold\textquotedblright{}.

\begin{quote}
\mr{निवेदकाने} \mr{आजच्या} \mr{ठळक} \mr{बातम्या} \mr{वाचल्या}.
\end{quote}

\begin{quote}
\emph{nivedakāne ājacyā ṭhaḷak bātamyā vācalyā.}
\end{quote}

\begin{quote}
The newsreader read out today's main headlines.
\end{quote}

It keeps one form: \mr{ठळक} \mr{अक्षरं}, bold letters; \mr{ठळक} \mr{बातम्या}, the headlines.

\begin{cousinweb}
\textbf{\mr{ठळक}} (\emph{ṭhaḷak}) is a Marathi word for what stands out clearly; its origin is not settled.
\end{cousinweb}

\begin{grammarlens}[title={What you've built}]
One more word for telling what happened.
\end{grammarlens}

\subsection*{Your turn}
\begin{itemize}
\raggedright
  \item \emph{Say it:} \emph{ṭhaḷak}
  \item \emph{Say it:} \emph{ṭhaḷak}, once more
  \item \emph{Say it:} \emph{nivedakāne ājacyā ṭhaḷak bātamyā vācalyā.} — The newsreader read out today's main headlines.
  \item \emph{Recall:} say the Marathi for a sensation, then the Marathi for an announcement, then say \emph{ṭhaḷak} again
\end{itemize}

\begin{tcolorbox}[breakable,colback=teal!4,colframe=teal!35!black,title={Before you move on}]
What does \mr{ठळक} mean? (\textquotedblleft{}Prominent; bold\textquotedblright{}.) Say it once more.
\end{tcolorbox}
\section[\texorpdfstring{\mr{पुकारणे} (pukārṇe)}{पुकारणे (pukārṇe)}]{\mr{पुकारणे} (pukārṇe) — to proclaim; to call a strike}

\label{lesson:MR-C313-pukarne}



\begin{quote}
Before the new one: say the Marathi for an announcement, then the Marathi for prominent.
\end{quote}

\subsection*{You'll want to know: \mr{पुकारणे}}
\textbf{\mr{पुकारणे}} — \emph{pukārṇe} — \textquotedblleft{}to proclaim; to call a strike\textquotedblright{}.

\begin{quote}
\mr{कामगारांनी} \mr{उद्यापासून} \mr{संप} \mr{पुकारला}.
\end{quote}

\begin{quote}
\emph{kāmgārānnī udyāpāsūn samp pukārlā.}
\end{quote}

\begin{quote}
The workers called a strike from tomorrow.
\end{quote}

Past, with -\mr{ने} (he …): \mr{त्याने} \mr{पुकारलं}, \emph{tyāne pukārlaṁ}.

\mr{संप} \mr{पुकारणे}, to call a strike; it also means to call out to someone.

\begin{cousinweb}
\textbf{\mr{पुकारणे}} (\emph{pukārṇe}) is inherited: CDIAL lists Marathi \emph{pukārṇe}, 'to shout', under \emph{*pukkār}, 'shout', a word made from the sound (CDIAL 8246).
\end{cousinweb}

\begin{grammarlens}[title={What you've built}]
One more word for telling what happened.
\end{grammarlens}

\subsection*{Your turn}
\begin{itemize}
\raggedright
  \item \emph{Say it:} \emph{pukārṇe}
  \item \emph{Say it:} \emph{pukārṇe}, once more
  \item \emph{Say it:} \emph{kāmgārānnī udyāpāsūn samp pukārlā.} — The workers called a strike from tomorrow.
  \item \emph{Recall:} say the Marathi for an announcement, then the Marathi for prominent, then say \emph{pukārṇe} again
\end{itemize}

\begin{tcolorbox}[breakable,colback=teal!4,colframe=teal!35!black,title={Before you move on}]
What does \mr{पुकारणे} mean? (\textquotedblleft{}To proclaim; to call a strike\textquotedblright{}.) Say it once more.
\end{tcolorbox}
